% Compile twice with: pdflatex Review_Report_Class_Example.tex
% Keep ssnreview.cls and the optional logo in the same folder.
\documentclass{ssnreview}
\usepackage{amsmath,amssymb}

\projecttitle{Reasoning Failure Boundaries in ECG\\
Foundation Models Under Progressive\\
Physiological Information Loss}
\reviewnumber{REVIEW REPORT -- II}
\students{%
  Vijay K (3122235001157)\\
  Siddharth Chandrasekar S (3122235001310)\\
  Vedanth M S (3122235001314)}
\supervisor{Dr. Anusha Jayasimhan\\Assistant Professor}
\department{Department of Computer Science and Engineering}
\institution{Sri Sivasubramaniya Nadar College of Engineering}
\academicyear{2026--2027}
\collegelogo{ssn-logo.png}

\begin{document}
\makereviewtitle
\reviewfrontmatter

\begin{reviewabstract}
ECG foundation models are increasingly used for automated interpretation and clinical
decision support, making their reliability under incomplete physiological observations an
important concern. In real-world acquisition and monitoring, ECG leads may be missing,
corrupted, or unavailable, yet a model may still produce confident interpretations from
insufficient evidence. Existing foundation and knowledge-enhanced approaches perform
strongly in representation learning, diagnosis, partial-lead analysis, and structured
interpretation, but are evaluated mainly through diagnostic accuracy and representation
quality. Consequently, how evidence-grounded reasoning degrades as physiological
information is removed remains insufficiently characterised.

This project investigates a knowledge-aware ECG reasoning framework under progressive
physiological information loss. Raw twelve-lead recordings are masked at seven availability
levels---12, 10, 8, 6, 4, 2 and 1---without reconstructing missing signals. Each degraded
recording is encoded by a frozen ECG foundation model into a latent representation,
converted into structured clinical findings, and passed to an evidence-based reasoning layer.
Matched conditions compare ECG-only reasoning with reasoning augmented by a Clinical
Knowledge Prior derived from SNOMED CT and ECG-specific clinical relationships.
The resulting diagnoses, supporting and opposing evidence, confidence, and uncertainty
are independently verified against an evidence schema and the lead-availability mask.
PTB-XL supports framework development, MIMIC-IV-ECG large-scale validation, and
Chapman-Shaoxing external validation.

Evaluation will measure grounding quality, evidence observability, diagnostic accuracy,
calibration, uncertainty, the Area Under the Reasoning Degradation Curve, and the
Reasoning Failure Boundary. The main contribution is a reproducible framework for
locating the boundary at which an ECG foundation model can no longer maintain
evidence-grounded reasoning and for testing whether structured clinical knowledge shifts that
boundary without promoting unsupported inference.

\keywords{Electrocardiogram; Foundation models; Evidence grounding; Reasoning
failure boundary; Clinical knowledge prior}
\end{reviewabstract}

\chapter{Introduction}

\subsection*{Project Domain and Context}

This project lies at the intersection of cardiovascular informatics, deep learning
foundation models, and machine-learning safety in healthcare. Its primary focus is the
automated interpretation of multi-lead electrocardiograms (ECGs) and, more importantly,
the verification of whether an artificial intelligence system's clinical reasoning remains
grounded in the physiological evidence that is actually available to it. The work therefore
extends beyond conventional diagnostic classification: it studies the reliability,
traceability, and safety of model-generated clinical claims under progressively degraded
and resource-constrained acquisition conditions.

The standard 12-lead ECG is one of the most widely used non-invasive tools for assessing
the electrical activity of the heart. The limb leads observe cardiac electrical activity
primarily in the frontal plane, whereas the precordial leads provide complementary views
across the horizontal plane. Together, these leads allow clinicians to examine cardiac
rhythm, conduction abnormalities, chamber-related changes, and regional patterns associated
with conditions such as myocardial ischaemia or infarction. Recent advances in deep learning
have enabled large-scale ECG foundation models (such as ECG-FM, K-MERL, and GEM) to learn
general-purpose representations from extensive collections of recordings, demonstrating
strong performance on rhythm classification and abnormality detection when complete,
high-quality recordings are available.

In practice, however, a complete and artifact-free 12-lead recording cannot always be
obtained. Emergency departments, pre-hospital transport services, intensive-care units,
ambulances, and wearable monitoring systems frequently operate under acquisition
constraints. Electrodes may become detached, signals may be corrupted by motion, and
time-critical triage may limit the number of leads that can be collected. Consumer and
portable devices also use reduced configurations consisting of six, four, two, or even a
single lead. These conditions create \emph{progressive physiological information loss}: as
leads are removed, the model loses not only signal samples but also specific anatomical
views of cardiac activity. A safe ECG reasoning system must recognise this change in
observability and adjust its claims, confidence, and uncertainty accordingly.

\subsection*{Motivation and Significance}

Most existing ECG artificial-intelligence benchmarks evaluate models using final-label
metrics such as accuracy, sensitivity, specificity, F1-score, or area under the receiver
operating characteristic curve. These measures assess predictive performance, but they do
not determine whether the reasoning used to reach a prediction is physically supported by the
active sensor configuration. 

This limitation becomes critical when the evidence required for a diagnosis is regional.
For example, confirmation of an inferior myocardial infarction conventionally depends on
observing compatible ST-segment changes across contiguous inferior leads, particularly
Leads II, III, and aVF. If key inferior leads such as III and aVF are unavailable, the
regional evidence needed to confirm the finding is no longer fully observable. A model may
nevertheless continue to output the correct diagnostic label because of correlations learned
from the training data, patient-level patterns, or dataset-specific shortcuts. It may also
generate a clinical explanation that refers to changes in leads that were not provided as
input. Conventional label-based evaluation mistakenly rewards such an output because the
predicted diagnosis matches the reference label. From a clinical-safety perspective,
however, a correct conclusion supported by unavailable or fabricated evidence is an
unacceptable silent failure: the output appears plausible and confident even though the
model has crossed the limit at which physical observations can justify its reasoning.

This project addresses this critical gap by formalising the \emph{Reasoning Failure
Boundary} (RFB), which identifies the level of lead availability beyond which the system can
no longer satisfy predefined grounding requirements. In addition, the \emph{Area Under the
Reasoning Degradation Curve} (AURDC) summarises how reasoning quality changes across the
complete sequence of lead-degradation conditions. The framework also introduces an explicit
three-state verification mechanism that labels each clinical claim as \textsc{Supported},
\textsc{Contradicted}, or \textsc{Unavailable}. This distinction prevents the absence of
evidence from being incorrectly interpreted as negative evidence. By incorporating
structured relationships derived from SNOMED CT through a Clinical Knowledge Prior (CKP),
the study evaluates whether clinical knowledge can help preserve coherent reasoning as sensor
information decreases, providing clinicians with interpretable evidence trails and
empirically measured operating limits.

\section{Objectives}
\begin{enumerate}[label=\arabic*.,leftmargin=*,align=left]
  \item To design and develop a controlled degradation framework and an architecture
  with two complementary representations that model progressive physiological information loss
  across standard ECG configurations containing 12, 10, 8, 6, 4, 2, and 1 available
  leads. The framework will remove leads without introducing synthetic or reconstructed
  signals and will combine continuous latent feature embeddings with discrete,
  deterministic clinical observations.

  \item To implement and evaluate a structured reasoning pipeline with a deterministic
  three-state grounding-verification mechanism that classifies individual clinical claims
  as \textsc{Supported}, \textsc{Contradicted}, or \textsc{Unavailable}. This mechanism
  will be used to determine whether the model's findings, evidence, and conclusions remain
  truthfully anchored in the physical sensor data at every degradation level.

  \item To compare an unaugmented ECG-only baseline with a knowledge-enhanced model that
  integrates a Clinical Knowledge Prior derived from the SNOMED CT clinical ontology into
  a shared reasoning transformer. The comparison will quantify changes in the Area Under
  the Reasoning Degradation Curve and determine whether the knowledge prior delays the
  Reasoning Failure Boundary under severe sensor scarcity without promoting unsupported
  clinical inference.
\end{enumerate}

\chapter{First Review: Queries and Responses}

The first project review raised important questions concerning the clinical validity of the
verification process, the necessity of the proposed dual-condition architecture, and the
practical value of the Reasoning Failure Boundary. This chapter records the questions and
provides the corresponding technical responses incorporated into the project methodology.

\section{Medical Reliability of the Evidence Schema}

\textbf{Query 1: How will you ensure that the Evidence Schema in Module 5 is medically
reliable?}

\textbf{Response:} The Evidence Schema is designed and validated with direct guidance from
a qualified medical professional possessing clinical expertise in ECG interpretation. Clinical
involvement is essential because Module~5 determines whether a model-generated claim is
physically supported by the electrodynamic evidence available in a particular lead
configuration. An incorrectly defined relationship between a clinical finding and its
required leads would produce a misleading grounding verdict regardless of software correctness.

For every admissible ECG claim, the schema specifies the clinical finding being assessed,
the waveform measurement required, the anatomical lead set from which that measurement
must be obtained, and the formal decision criterion used to classify the claim. Particular
attention is given to regional findings, such as ST-segment changes, whose interpretation
depends on the presence of specific contiguous leads. 

The schema-development protocol incorporates four rigorous stages: (1) constructing an
initial machine-readable schema using established ECG terminology, SNOMED CT concepts, and
cardiology guidelines; (2) reviewing each claim type, required lead set, measurement
definition, and verification criterion with the medical expert; (3) validating the schema on
representative complete-lead and reduced-lead benchmark ECG recordings to verify that
\textsc{Supported}, \textsc{Contradicted}, and \textsc{Unavailable} verdicts are assigned
accurately; and (4) version-controlling all approved schema rules prior to model evaluation to
prevent post-hoc modifications.

\section{Need for the Control and Experimental Conditions}

\textbf{Query 2: Why is the model complex? Do we actually need both the control and
experimental parts?}

\textbf{Response:} Both reasoning conditions are strictly necessary to maintain experimental
control while investigating whether structured knowledge can mitigate information loss.
The ECG-only control path establishes how the foundation model's reasoning degrades purely
as physiological information is removed, producing the reference degradation curve and
baseline Reasoning Failure Boundary. 

The experimental CKP path addresses a distinct research question: whether structured
clinical knowledge can help preserve valid reasoning under reduced-lead conditions without
encouraging unsupported inference. The Clinical Knowledge Prior does not reconstruct missing
leads or supply synthetic waveform observations; instead, it provides structured relationships
among diagnoses, clinical findings, and anatomical dependencies that guide the reasoning
component in organizing the evidence still available to it.

Although the architecture contains several modules, this separation enforces experimental
auditability. Signal preprocessing, controlled masking, latent representation, observation
extraction, reasoning, and grounding verification are isolated so that failures can be
localized to a specific stage. Crucially, the two reasoning conditions share the same patient
recordings, masks, frozen ECG encoder, latent representations, structured observations,
candidate diagnoses, transformer architecture, and verification rules. They differ exclusively
through the inclusion of CKP relational tokens in the experimental context, ensuring that any
measured difference in grounding performance is attributable strictly to the knowledge prior
rather than confounding model capacity.

\section{Purpose and Clinical Utility of the Reasoning Failure Boundary}

\textbf{Query 3: What are you achieving with this work? How will the Reasoning Failure
Boundary be useful?}

\textbf{Response:} The project aims to determine how much observable ECG information is
required for a clinical AI system to produce reasoning that remains supported by physical
sensor data, distinguishing diagnostic correctness from evidence-grounded reliability.

The Reasoning Failure Boundary (RFB) provides an empirical operating limit. For a
predefined clinical-safety threshold (e.g., Grounding Quality $\ge 0.60$), it identifies the
most severely degraded lead configuration at which the system's reasoning remains acceptable.
For example, if the model meets the grounding threshold with six available leads but fails at
four leads, the boundary indicates that the model should not be deployed below the six-lead
configuration. The complete reasoning degradation curve and AURDC are reported alongside the
boundary so that conclusions do not depend on a single operating point.

This boundary assists clinicians and developers in several critical ways: (1) it indicates
when automated systems in ambulances or bedside telemetry must alert clinicians to acquire a
complete 12-lead ECG; (2) it identifies when the system should reduce its confidence or
abstain from diagnosis; and (3) it defines certified operational envelopes for wearable and
low-channel monitoring devices. The RFB converts an otherwise hidden failure mode into an
interpretable, measurable safety indicator.

\chapter{Problem statement}

Modern ECG foundation models can maintain apparently strong diagnostic performance even
when a substantial portion of the original sensor information is unavailable. Under
progressive lead loss, a model may continue to predict the correct diagnostic label by
exploiting statistical regularities, dataset-specific shortcuts, or correlations that do not
constitute direct physiological evidence for the prediction. More critically, a generative
reasoning component may produce clinical claims referring to waveform changes in leads
that were not supplied to the model. Conventional classification metrics do not distinguish
between a diagnosis supported by observable electrodynamic evidence and a plausible guess
supported by an unavailable or hallucinated finding.

This limitation represents a significant safety concern in emergency triage, pre-hospital
transport, intensive-care monitoring, and wearable telemetry, where missing or detached
electrodes are common. When the required anatomical projections are absent, the available
recording may no longer contain sufficient information to confirm or reject a localised
pathology. A reliable clinical system must therefore recognise not only what diagnosis is
probable, but also whether the evidence needed to justify that diagnosis can be observed
from the active physical sensors.

The problem addressed in this project is the absence of a systematic method for measuring
how evidence-grounded clinical reasoning degrades as ECG leads are progressively removed.
The project seeks to determine the operating point at which a foundation model's reasoning
can no longer be considered sufficiently grounded, even if its final diagnostic accuracy
remains high. It also investigates whether structured clinical knowledge can delay this
failure point without replacing missing signals with artificial evidence or encouraging
unsupported inference.

Formally, the study considers the ordered set of lead-availability levels:
\begin{equation}
  \mathcal{L} = (12, 10, 8, 6, 4, 2, 1),
\end{equation}
where each value denotes the number of physical ECG leads retained at a particular
degradation stage. The same patient recording is evaluated across these controlled stages,
allowing diagnostic performance, evidence observability, grounding quality, and uncertainty
to be measured as functions of the available physiological information under two matched
reasoning conditions: an ECG-only baseline and a knowledge-guided condition augmented with a
Clinical Knowledge Prior (CKP) derived from SNOMED CT.

\section{Inputs}

The proposed system receives four principal categories of input:

\paragraph{Multi-lead ECG Waveform:} A multi-channel surface ECG recording represented as a
voltage-time matrix. Every recording is standardised to a fixed 10.0-second duration and
500~Hz sampling frequency ($12 \times 5000$ samples) so that all degradation conditions are
evaluated over the same temporal interval. Signal values are retained only for leads that are
physically available at the selected degradation level; unavailable leads are not
regenerated, inferred, or replaced with synthetic waveforms. Retained signals undergo a
consistent preprocessing pipeline for quality control, amplitude normalisation,
baseline-wander removal, and suppression of high-frequency interference while preserving
diagnostically vital P-wave, QRS-complex, ST-segment, and T-wave morphology.

\paragraph{Lead-Availability Mask:} Each recording is accompanied by a binary lead-availability
mask:
\begin{equation}
  \mathbf{m}^{(\ell)} = (m_1,m_2,\ldots,m_{12}), \qquad
  m_j \in \{0,1\}, \quad \ell \in \mathcal{L},
\end{equation}
where $m_j=1$ indicates that lead $j$ is physically available and $m_j=0$ indicates that it
has been removed. The mask represents physical sensor observability rather than a request to
impute missing channels. It is propagated through all reasoning and verification stages so that
every clinical claim can be audited against the anatomical views present in the input.

\paragraph{Candidate Diagnostic Hypotheses:} A predefined candidate set covering major
cardiac disease and rhythm categories represented in the evaluation datasets. These candidates
define the shared diagnostic space over which posterior probabilities are estimated for both
conditions, preventing vocabulary discrepancies from influencing comparative evaluation.

\paragraph{Clinical Knowledge Prior:} A structured ontological representation derived from
SNOMED CT and ECG-specific clinical relationships, encoding associations among diagnoses,
observable findings, anatomical regions, and verification rules. Its role is to constrain and
organise reasoning; it is strictly prohibited from synthesising missing signals or declaring
unobserved findings.

\section{Outputs}

For every recording and lead-availability level, the system produces four integrated outputs:

\paragraph{Diagnostic Posterior Distribution:} A ranked posterior probability distribution
over candidate cardiac diagnoses, allowing assessment of diagnostic discrimination,
calibration, and changes in confidence or diagnostic abstention as leads are removed.

\paragraph{Structured Clinical Reasoning Trace:} An interpretable reasoning trace composed of
discrete clinical claims detailing physiological observations, associated candidate
diagnoses, supporting or opposing status, contributing leads, and cross-attention weights.

\paragraph{Grounding-Verification Verdicts:} A deterministic verifier assigns every generated
claim to exactly one of three states:
\begin{itemize}
  \item \textsc{Supported}: The required leads are available and the claimed physiological
  pattern is verified in the observed signal;
  \item \textsc{Contradicted}: The required leads are available, but the observed waveform
  conflicts with the generated claim; or
  \item \textsc{Unavailable}: The claim cannot be evaluated because one or more anatomical
  projections required to verify it are absent.
\end{itemize}
This three-state taxonomy prevents unavailable evidence from being misinterpreted as
negative evidence and detects when the model reasons beyond its visible sensor space.

\paragraph{Reasoning-Degradation Metrics:} The claim-level verdicts are aggregated at each lead
level to compute: \emph{Evidence Observability} (proportion of diagnostic evidence assessable
under the active lead set), \emph{Grounding Quality} (precision of reasoning against
observable evidence), \emph{True Grounding Ratio} (proportion of claims explicitly supported
by retained signals), \emph{Area Under the Reasoning Degradation Curve} (normalised summary
across all seven levels), and the \emph{Reasoning Failure Boundary} (the first lead level where
Grounding Quality falls below the clinical safety threshold).

\section{Constraints}

The investigation is conducted under five rigorous operational and scientific constraints:
\begin{itemize}
  \item \textbf{No Reconstruction of Unavailable Leads:} Generative lead reconstruction, deep
  inpainting, mathematical lead synthesis, and substitution using predicted waveforms are
  strictly prohibited. Missing leads remain missing to benchmark authentic information scarcity.
  \item \textbf{Morphology-Preserving Preprocessing:} Digital filters are applied in a
  zero-phase, bidirectional manner (\texttt{scipy.signal.filtfilt}) to eliminate net phase delay
  and prevent artificial displacement of QRS boundaries or ST segments.
  \item \textbf{Patient-Stratified Evaluation:} Training, validation, and testing partitions are
  created strictly at the patient level, preventing identical-patient heart geometry leakage.
  \item \textbf{Controlled Experimental Comparison:} Both conditions share identical waveforms,
  lead masks, candidate diagnoses, preprocessing operations, and verification rules.
  \item \textbf{Evidence-Limited Use of Prior Knowledge:} The CKP guides the organisation of
  evidence but cannot override the lead-availability mask or convert an unobservable finding
  into a supported claim.
\end{itemize}

\section{Success Criteria}

The proposed system is evaluated using explicit and measurable safety criteria. Let
$GQ(\ell)$ denote Grounding Quality with $\ell$ available leads, and let the clinical safety
threshold be:
\begin{equation}
  \tau_{\mathrm{safety}} = 0.60.
\end{equation}
For lead levels $\mathcal{L}=\{12,10,8,6,4,2,1\}$, the Reasoning Failure Boundary is:
\begin{equation}
  \mathrm{RFB}=\min\{\ell\in\mathcal{L}:GQ(\ell)\geq 0.60\}.
\end{equation}
The knowledge-guided system is considered successful when it satisfies two primary criteria:
\begin{enumerate}[label=\arabic*.,leftmargin=*,align=left]
  \item It achieves a normalised AURDC of at least $0.80$ across the full sequence from twelve
  leads down to a single lead, demonstrating sustained reasoning fidelity.
  \item Its RFB reduces the minimum acceptable lead requirement by at least two physical leads
  relative to the unaugmented baseline (e.g., maintaining acceptable grounding at 2 leads if
  the baseline collapses at 4 leads).
\end{enumerate}
Secondary evaluation examines whether this improvement is achieved without compromising
diagnostic discrimination or calibration across \textsc{Supported}, \textsc{Contradicted},
and \textsc{Unavailable} claim states.

\chapter{Methodology}

The proposed framework is organised into six sequential, auditable phases. An ECG recording
is ingested and standardised, progressively degraded using controlled lead masking, encoded
into a latent representation by a frozen foundation model, processed under two matched
reasoning conditions, verified claim-by-claim against observable physical signals, and
summarised using degradation indices.

\section{Architectural Design for Proposed System}

Figure~\ref{fig:architecture} presents the high-level architecture of the proposed system,
linking data preparation, controlled masking, latent representation, clinical reasoning,
grounding verification, and results generation. Crucially, the architecture exposes a direct
data path from the masking phase to the grounding verification phase, ensuring the verifier
receives the ground-truth mask and physical signals directly rather than relying on model
self-reports.

\begin{figure}[H]
  \centering
  \includegraphics[width=\textwidth]{../Gemini_Generated_Image_e5sdgje5sdgje5sd.png}
  \caption{High-level architecture of the proposed system}
  \label{fig:architecture}
\end{figure}

\subsection*{End-to-End Evaluation Procedure}

For every ECG recording, the complete evaluation executes eight sequential steps:
\begin{enumerate}[label=\textbf{Step \arabic*:},leftmargin=*,align=left]
  \item Standardise the raw recording using the ordered preprocessing pipeline, isolating
  diagnostic labels as reference ground truth.
  \item For each lead level $\ell \in \Lambda = (12, 10, 8, 6, 4, 2, 1)$, generate a masked
  recording $\mathbf{X}_\ell$ and its binary availability mask $\mathbf{a}_\ell$.
  \item Encode the masked recording through the frozen ECG foundation model to obtain continuous
  latent representation $\mathbf{Z}_\ell$.
  \item Extract structured ECG observations $\mathcal{O}_\ell$ and construct a shared candidate
  hypothesis set $\mathcal{H}_\ell$.
  \item Generate one reasoning trace using the ECG-only control context and a second trace
  using the CKP-assisted experimental context.
  \item Decompose each trace into individual claims and verify each against the evidence
  schema, masked waveform, and availability mask.
  \item Aggregate claim verdicts into Evidence Observability ($EO_\ell$) and Grounding Quality
  ($GQ_\ell$) for each lead level and condition.
  \item Construct the reasoning degradation curves, compute AURDC and RFB, and quantify the
  paired improvement produced by the knowledge prior.
\end{enumerate}

\section{Modules}

Let $\mathcal{C} = \{\mathrm{I},\mathrm{II},\mathrm{III},\mathrm{aVR},\mathrm{aVL},\mathrm{aVF},\mathrm{V1},\ldots,\mathrm{V6}\}$
denote the standard 12-lead set, and $\Lambda = (12, 10, 8, 6, 4, 2, 1)$ denote the evaluated
lead levels. A standardised recording is represented by $\mathbf{X}\in\mathbb{R}^{12\times n_T}$.
At level $\ell\in\Lambda$, $\mathbf{a}_\ell\in\{0,1\}^{12}$ is the availability mask, and:
\begin{equation}
  \mathcal{A}_\ell = \{\lambda\in\mathcal{C}:\mathbf{a}_\ell[\lambda]=1\}
\end{equation}
is the active physical lead set.

\subsection{Module 1: Data and Input Management}
\textbf{Purpose \& Processing:} Establishes a reproducible data pipeline strictly separating
observed signals from evaluation annotations. Five ordered operations standardize heterogeneous
recordings from PTB-XL, MIMIC-IV-ECG, and Chapman-Shaoxing:
\begin{equation}
  \mathbf{X} = (P_5 \circ P_4 \circ P_3 \circ P_2 \circ P_1)(\mathbf{X}_{\mathrm{raw}}),
\end{equation}
where operators enforce canonical lead ordering ($P_1$), sampling and duration standardization
($P_2$), voltage calibration and normalization ($P_3$), zero-phase digital filtering ($P_4$),
and quality screening ($P_5$). Outputs standardized tensors of shape $12 \times n_T$ and
reproducibility manifests.

\subsection{Module 2: Controlled Degradation}
\textbf{Purpose \& Processing:} Defines the primary independent variable by systematically
masking leads across the seven availability levels. Inactive channels are clamped to a fixed
constant representation $\nu = 0$ while the mask $\mathbf{a}_\ell$ is retained as an explicit
input, avoiding synthetic reconstruction:
\begin{equation}
  \mathbf{X}_\ell = \mathcal{M}(\mathbf{X}, \mathbf{a}_\ell) = \mathrm{diag}(\mathbf{a}_\ell)\mathbf{X} + \nu(\mathbf{1}_{12} - \mathbf{a}_\ell)\mathbf{1}_{n_T}^\top, \quad |\mathcal{A}_\ell| = \ell.
\end{equation}
Outputs seven degraded variants $\{(\mathbf{X}_\ell, \mathbf{a}_\ell) : \ell \in \Lambda\}$ per recording.

\subsection{Module 3: Latent Representation}
\textbf{Purpose \& Processing:} Converts degraded waveforms into continuous latent vectors using
frozen foundation model weights (ECG-FM), preventing representation adaptation from
confounding reasoning degradation:
\begin{equation}
  \mathbf{Z}_\ell = f_\theta(\mathbf{X}_\ell, \mathbf{a}_\ell) \in \mathbb{R}^D, \quad \theta \text{ fixed}.
\end{equation}
Masking sensitivity is audited via cosine divergence from the full 12-lead representation:
\begin{equation}
  \Delta_\ell^{\cos} = 1 - \frac{\langle\mathbf{Z}_{12},\mathbf{Z}_\ell\rangle}{\|\mathbf{Z}_{12}\|_2\|\mathbf{Z}_\ell\|_2}.
\end{equation}

\subsection{Module 4: Clinical Reasoning}
\textbf{Purpose \& Processing:} Projects the latent vector into discrete clinical observation
fields $\mathcal{F} = \{\text{rhythm}, \text{rate}, \text{PR}, \text{QRS}, \text{ST}\}$:
\begin{equation}
  \mathbf{z}_\ell^{(f)} = \mathbf{W}_f\mathbf{Z}_\ell + \mathbf{b}_f, \qquad
  \hat{y}_\ell^{(f)} = \arg\max \mathrm{softmax}(\mathbf{z}_\ell^{(f)}).
\end{equation}
A shared rule base $\mathcal{R}$ converts observations $\mathcal{O}_\ell$ into candidate
diagnoses $\mathcal{H}_\ell = \mathcal{R}(\mathcal{O}_\ell) = \{h_1,\dots,h_K\}$. For candidate
$h_i$, the relevant SNOMED CT subgraph $\mathcal{G}_{\ell,i}$ is encoded and additively fused:
\begin{equation}
  \mathbf{Z}_{\ell,i}^K = \mathbf{W}_K\phi(\mathcal{G}_{\ell,i}) + \mathbf{b}_K \in \mathbb{R}^D, \qquad
  \mathrm{CKP}_{\ell,i} = \mathbf{Z}_\ell + \mathbf{Z}_{\ell,i}^K.
\end{equation}
Matched reasoning contexts $\mathcal{C}_{\ell,i}^{(0)} = [\mathcal{O}_\ell; \mathcal{E}_{\ell,i}; h_i]$
and $\mathcal{C}_{\ell,i}^{(1)} = [\mathcal{O}_\ell; \mathrm{CKP}_{\ell,i}; \mathcal{E}_{\ell,i}; h_i]$
are evaluated by a shared transformer, yielding candidate scores and calibrated posterior
probabilities via temperature scaling:
\begin{equation}
  \mathbf{U}_{\ell,i}^{(c)} = \mathrm{Trf}_\varphi(\mathcal{C}_{\ell,i}^{(c)}), \quad
  s_{\ell,i}^{(c)} = \mathbf{w}_s^\top \mathbf{U}_{\ell,i}^{(c)} + b_s, \quad
  p_{\ell,i}^{(c)} = \frac{\exp(s_{\ell,i}^{(c)}/\beta_\ell)}{\sum_{i'=1}^K \exp(s_{\ell,i'}^{(c)}/\beta_\ell)}, \quad
  \hat{d}_\ell^{(c)} = \arg\max_{h_i\in\mathcal{H}_\ell} p_{\ell,i}^{(c)}.
\end{equation}

\subsection{Module 5: Grounding Verification}
\textbf{Purpose \& Processing:} Decomposes reasoning traces into individual claims $\gamma$ and
evaluates each against an independently defined evidence schema
$\mathcal{S}(\gamma) = (\mathrm{meas}(\gamma), \mathcal{R}(\gamma), \pi(\gamma))$, where
$\mathcal{R}(\gamma) \subseteq \mathcal{C}$ is the required lead set and $\pi(\gamma)$ is the
waveform criterion. The fraction of observable evidence is:
\begin{equation}
  o(\gamma, \ell) = \frac{|\mathcal{R}(\gamma)\cap\mathcal{A}_\ell|}{|\mathcal{R}(\gamma)|} \in [0,1].
\end{equation}
If $o(\gamma,\ell) = 0$, the claim is classified as \textsc{Unavailable}. If $o(\gamma,\ell) > 0$,
the waveform is evaluated against $\pi(\gamma)$ to yield \textsc{Supported} or
\textsc{Contradicted}. For claim set $\Gamma_\ell$, required and observable evidence counts are:
\begin{equation}
  N_\ell^{\mathrm{req}} = \sum_{\gamma\in\Gamma_\ell} |\mathcal{R}(\gamma)|, \qquad
  N_\ell^{\mathrm{obs}} = \sum_{\gamma\in\Gamma_\ell} |\mathcal{R}(\gamma)\cap\mathcal{A}_\ell|.
\end{equation}
Level-wise Grounding Quality ($GQ_\ell$) and Evidence Observability ($EO_\ell$) are computed as:
\begin{equation}
  GQ_\ell = \frac{N_\ell^{\mathrm{sup}}}{N_\ell^{\mathrm{sup}} + N_\ell^{\mathrm{con}}}, \qquad
  EO_\ell = \frac{N_\ell^{\mathrm{obs}}}{N_\ell^{\mathrm{req}}}.
\end{equation}

\subsection{Module 6: Results and Evaluation}
\textbf{Purpose \& Processing:} Compiles per-level grounding scores into the reasoning
degradation curve $\mathcal{D}(\ell) = GQ_\ell$. Normalized AURDC across
$(L_1,\dots,L_7) = (12,10,8,6,4,2,1)$ is calculated using weighted trapezoidal integration:
\begin{equation}
  \mathrm{AURDC} = \frac{1}{11} \sum_{i=1}^6 \frac{GQ_{L_i} + GQ_{L_{i+1}}}{2}(L_i - L_{i+1}).
\end{equation}
The Reasoning Failure Boundary is identified as $\mathrm{RFB} = \min\{\ell\in\Lambda : GQ_\ell \ge \tau\}$,
and paired knowledge improvements are reported as:
\begin{equation}
  \Delta GQ_\ell = GQ_\ell^{\mathrm{CKP}} - GQ_\ell^{\mathrm{Control}}, \quad
  \Delta\mathrm{AURDC} = \mathrm{AURDC}^{\mathrm{CKP}} - \mathrm{AURDC}^{\mathrm{Control}}, \quad
  \Delta\mathrm{RFB}_{\mathrm{imp}} = \mathrm{RFB}^{\mathrm{Control}} - \mathrm{RFB}^{\mathrm{CKP}}.
\end{equation}

\chapter{Implementation}

\section{Module 1: Data Acquisition, Preprocessing, and Controlled Degradation}

Module~1 implements the data-engineering and signal-processing pipeline in Python~3.11 using
NumPy, SciPy, Pandas, and the Waveform Database (WFDB) library. Every transformation is recorded
in per-record manifests ensuring complete experimental auditability.

\subsection{Dataset Ingestion and Patient-Stratified Partitioning}

The primary development benchmark is PTB-XL version~1.0.3, containing 21,799 clinical 10-second
12-lead ECG recordings from 18,869 unique patients. Diagnostic statements follow the SCP-ECG
standard across five superclasses: Normal ECG (\textsc{NORM}), Myocardial Infarction
(\textsc{MI}), ST/T Changes (\textsc{STTC}), Conduction Disturbance (\textsc{CD}), and
Hypertrophy (\textsc{HYP}). Partitioning strictly enforces patient-level quarantine so that no
patient appears in multiple splits (Table~\ref{tab:ptbxl-splits}).

\begin{table}[H]
  \centering
  \small
  \caption{Patient-stratified PTB-XL partitions used by the implementation}
  \label{tab:ptbxl-splits}
  \begin{tabularx}{\textwidth}{@{}lccc>{\raggedright\arraybackslash}X@{}}
    \toprule
    \textbf{Partition} & \textbf{Folds} & \textbf{Records} &
    \textbf{Patients} & \textbf{Purpose} \\
    \midrule
    Training & 1--8 & 17,418 & 15,096 &
    Adapter optimisation and representation learning \\
    Validation & 9 & 2,183 & 1,893 &
    Hyperparameter selection, calibration, and early stopping \\
    Held-out test & 10 & 2,198 & 1,880 &
    Quarantined reasoning-degradation benchmark \\
    \bottomrule
  \end{tabularx}
\end{table}

\subsection{Five Sequential Preprocessing Operations}

The \texttt{ECGPreprocessingPipeline} executes five operations in strict sequence:

\paragraph{Operation P1: Lead Ordering and Spatial Alignment:} Maps source channel labels to the
canonical sequence:
\begin{equation}
  (\mathrm{I},\mathrm{II},\mathrm{III},\mathrm{aVR},\mathrm{aVL},
  \mathrm{aVF},\mathrm{V1},\mathrm{V2},\mathrm{V3},\mathrm{V4},
  \mathrm{V5},\mathrm{V6}),
\end{equation}
transposing arrays into a channel-first layout of shape $(12, T)$ and raising a strict
\texttt{LeadOrderingError} upon channel absence or duplication.

\paragraph{Operation P2: Sampling-Rate and Duration Standardisation:} Converts heterogeneous
input rates to 500~Hz using zero-phase Fourier-domain interpolation
(\texttt{scipy.signal.resample}). Centered temporal windowing crops signals $>10.0$~s around
their midpoint, while shorter signals are symmetrically zero-padded, yielding fixed dimensions:
\begin{equation}
  12\times(500\ \mathrm{samples/s}\times10\ \mathrm{s}) = 12\times5000.
\end{equation}

\paragraph{Operation P3: Unit Calibration and Physical-Voltage Normalisation:} Converts raw ADC
counts into physical millivolts and generates standardized neural inputs:
\begin{equation}
  V_{\mathrm{mV}}[n] = \frac{x_{\mathrm{raw}}[n]-b_{\mathrm{ADC}}}{g_{\mathrm{ADC}}}, \qquad
  \widetilde{V}_{c}[n] = \frac{V_c[n]-\mu_c}{\sigma_c+\varepsilon},
  \label{eq:adc-calibration}
\end{equation}
with physiological outlier clamping to $[-15.0, +15.0]$~mV.

\paragraph{Operation P4: Zero-Phase Digital Filtering:} Applies bidirectional filtering
(\texttt{scipy.signal.filtfilt}) to prevent group delay and phase distortion: a 3rd-order
Butterworth high-pass filter at 0.5~Hz (baseline wander), 2nd-order IIR notch filters at
50/60~Hz ($Q=30$, mains interference), and a 4th-order Butterworth low-pass filter at 100~Hz
(EMG noise).

\paragraph{Operation P5: Quality Screening and Integrity Audit:} Four automated gates validate
the signal: dimensional shape ($12 \times 5000$), non-finite rejection (detecting
$\mathrm{NaN}/\pm\infty$), flatline detection ($\sigma_c < 10^{-4}$~mV or static run $>1.5$~s),
and amplifier saturation clipping ($|V| \ge 15$~mV). Emits a structured \texttt{QualityAuditResult}.

\subsection{Controlled Lead-Availability Masking}

Clean recordings pass to the degradation engine, which applies binary masks across
$\Lambda = (12, 10, 8, 6, 4, 2, 1)$:
\begin{equation}
  \widetilde{\mathbf{X}}_\ell[c, n] = \mathbf{a}_\ell[c] \widetilde{\mathbf{X}}[c, n].
\end{equation}
Inactive channels are zero-clamped without generative inpainting or reconstruction, preserving
genuine sensor scarcity.

\section{Module 2: Core Reasoning Engine}

\subsection{Foundation-Model Backbone and Residual Bottleneck Adapter}

The architecture supports a 1D-ResNet convolutional baseline alongside the pre-trained ECG-FM
foundation model (12-layer Transformer encoder, 16 attention heads, 768 hidden dimension).
A lightweight 2-layer residual bottleneck MLP adapter (Table~\ref{tab:adapter-config}) projects
backbone activations into the standardized 768-dimensional latent vector
$\mathbf{Z}_{\mathrm{ECG}} \in \mathbb{R}^{768}$.

\begin{table}[H]
  \centering
  \small
  \caption{Residual adapter optimisation configuration}
  \label{tab:adapter-config}
  \begin{tabular}{@{}ll@{}}
    \toprule
    \textbf{Parameter} & \textbf{Configuration} \\
    \midrule
    Training records & 17,418 recordings from Folds 1--8 \\
    Validation records & 2,183 recordings from Fold 9 \\
    Optimiser & AdamW \\
    Learning rate & $1\times 10^{-4}$ \\
    Weight decay & $1\times 10^{-2}$ \\
    Maximum epochs & 20 \\
    Dropout probability & 0.20 \\
    Objective & Multi-label binary cross-entropy with logits \\
    Output dimension & 768 \\
    \bottomrule
  \end{tabular}
\end{table}

The adapter achieved a validation loss of 0.3057 and validation Macro-AUC of 0.7628 across the
8 benchmark classes. Latent representations were extracted for all 2,198 test records across
all degradation levels, producing numerical matrices of shape:
\begin{equation}
  \mathbf{Z}_{\mathrm{test}}^{(\ell)} \in \mathbb{R}^{2198\times 768}
\end{equation}
with verified numerical stability and zero non-finite values.

\subsection{Deterministic Clinical-Observation Extraction}

The \texttt{ClinicalObservationExtractor} extracts discrete biomarkers strictly from unmasked
leads:
\begin{itemize}
  \item \textbf{Cardiac Rhythm and Rate:} Evaluates R-peaks from physical-unit signals and computes
  heart rate via median R--R intervals:
  \begin{equation}
    \mathrm{HR} = \frac{60f_s}{\mathrm{median}(\Delta n_{\mathrm{RR}})}\ \mathrm{beats/min}.
  \end{equation}
  \item \textbf{Conduction Intervals:} Measures PR interval and QRS duration in milliseconds.
  \item \textbf{ST-Segment and T-Wave Morphology:} Calculates ST displacement at the J-point plus
  80~ms relative to the isoelectric TP baseline:
  \begin{equation}
    \Delta V_{\mathrm{ST},c} = V_c(t_J + 80\ \mathrm{ms}) - V_{\mathrm{TP},c},
  \end{equation}
  recording findings as \textsc{Unavailable} whenever required anatomical channels are absent.
\end{itemize}

\subsection{Ontological Prior and Dual-Path Reasoning Engine}

SNOMED CT relations formalize disease rules (e.g., Inferior MI requiring ST elevation $\ge 1.0$~mm
in $\ge 2$ contiguous leads: II, III, aVF). The shared transformer evaluates candidate
hypotheses under the ECG-only control context and the CKP-augmented experimental context,
emitting class posterior probabilities and evidence cross-attention matrices.

\section{Module 3: User Interface, Telemetry, and Output Management}

\subsection{Structured Clinical-Reasoning Trace}

Module~3 produces interpretable diagnostic traces separating diagnostic hypotheses from underlying
evidence (Table~\ref{tab:trace-fields}).

\begin{table}[H]
  \centering
  \small
  \caption{Principal fields in the structured clinical-reasoning trace}
  \label{tab:trace-fields}
  \begin{tabularx}{\textwidth}{@{}l>{\raggedright\arraybackslash}X@{}}
    \toprule
    \textbf{Field} & \textbf{Description} \\
    \midrule
    Record context & Record identifier, patient-safe identifier, split, lead level,
    active leads, model checkpoint, and reasoning condition \\
    Primary hypothesis & Highest-ranked cardiac diagnosis and calibrated posterior probability \\
    Alternatives & Remaining candidate diagnoses and their probabilities \\
    Supporting evidence & Structured physiological claims supporting the hypothesis,
    with contributing leads and attention weights \\
    Opposing evidence & Findings that reduce support for the hypothesis \\
    Missing evidence & Required clinical observations that cannot be evaluated from the active sensors \\
    Uncertainty & Predictive uncertainty and information-scarcity indicators \\
    Verification & Claim-level \textsc{Supported}, \textsc{Contradicted}, or \textsc{Unavailable} verdicts \\
    \bottomrule
  \end{tabularx}
\end{table}

\subsection{Telemetry, Reporting, and Visualisation}

The telemetry layer passes claims to the three-state grounding gate, reporting:
\begin{equation}
  \text{Grounded Support Ratio} = \frac{N_{\mathrm{sup}}}{N_{\mathrm{all}}}, \quad
  \text{Contradicted Ratio} = \frac{N_{\mathrm{con}}}{N_{\mathrm{all}}}, \quad
  \text{Unavailable Ratio} = \frac{N_{\mathrm{unav}}}{N_{\mathrm{all}}}.
\end{equation}
The reporting engine serializes results into cohort inventories, JSON traces, and CSV summary
tables. The plotting subsystem (Matplotlib/Seaborn) generates cohort distribution figures,
representation diagnostics, and Reasoning Degradation Curves plotting Grounding Quality against
available leads.

\chapter{Results and Discussion}

\section{Experimental Setup and Evaluation Protocol}

The proposed framework was evaluated using the held-out Fold~10 partition of PTB-XL,
comprising 2,198 clinical recordings from 1,880 unique patients \cite{ptbxl2020}.
Patient-stratified isolation was preserved throughout training, validation, calibration, and
testing. Every test recording was evaluated at seven progressively degraded lead configurations:
12, 10, 8, 6, 4, 2, and 1 lead, comparing the ECG-only control against the CKP-augmented
experimental condition.

Evaluation jointly measures correctness and observability: Grounding Quality (GQ), Evidence
Observability (EO), True Grounding Ratio:
\begin{equation}
  \mathrm{TGR}_\ell = \mathrm{GQ}_\ell \times \mathrm{EO}_\ell,
  \label{eq:tgr-results}
\end{equation}
Area Under the Reasoning Degradation Curve (AURDC), and the Reasoning Failure Boundary (RFB at
$\tau_{\mathrm{safety}}=0.60$).

\section{Quantitative Results and Degradation Analysis}

\subsection{Evidence Observability and Physical Information Loss}

Table~\ref{tab:evidence-retention} summarizes observation-extraction results across the seven
lead configurations.

\begin{table}[H]
  \centering
  \small
  \caption{Observable ST evidence under progressive lead degradation}
  \label{tab:evidence-retention}
  \begin{tabularx}{\textwidth}{@{}c>{\raggedright\arraybackslash}Xrrr@{}}
    \toprule
    \textbf{Leads} & \textbf{Active configuration} &
    \textbf{ST elevation} & \textbf{ST depression} &
    \textbf{EO} \\
    \midrule
    12 & Complete standard 12-lead ECG & 53.69\% & 82.30\% & 100.0\% \\
    10 & Leads III and aVF masked & 52.50\% & 82.07\% & 97.8\% \\
    8 & Derived limb leads removed & 37.58\% & 80.85\% & 70.0\% \\
    6 & I, II, V1, V2, V5, and V6 & 36.35\% & 80.35\% & 67.7\% \\
    4 & I, II, V1, and V5 & 25.89\% & 79.39\% & 48.2\% \\
    2 & Telemetry configuration: II and V1 & 22.47\% & 62.60\% & 41.9\% \\
    1 & Ambulatory configuration: II & 0.86\% & 62.06\% & 1.6\% \\
    \bottomrule
  \end{tabularx}
\end{table}

Lead removal induces an acute loss of anatomical evidence rather than a mere reduction in sample
count. Evidence retention declines from 100\% at 12 leads to 41.9\% at 2 leads, before collapsing
to 1.6\% in the single-lead condition. Crucially, ST elevation evidence collapses from 53.69\%
down to 0.86\% with Lead~II alone. Because regional myocardial injury requires spatial agreement
across contiguous leads, single-lead monitoring physically eliminates the physiological evidence
needed to substantiate regional infarction.

\subsection{Grounding Quality and True Grounding}

Table~\ref{tab:grounding-results} compares the two reasoning conditions across degradation levels.

\begin{table}[H]
  \centering
  \small
  \caption{Grounding degradation of the control and CKP-assisted conditions}
  \label{tab:grounding-results}
  \begin{tabular}{@{}crrrrr@{}}
    \toprule
    \textbf{Leads} &
    \textbf{GQ Control} &
    \textbf{GQ CKP} &
    \textbf{Margin} &
    \textbf{TGR Control} &
    \textbf{TGR CKP} \\
    \midrule
    12 & 0.88 & 0.91 & +0.03 & 0.880 & 0.910 \\
    10 & 0.80 & 0.86 & +0.06 & 0.782 & 0.841 \\
    8  & 0.72 & 0.77 & +0.05 & 0.504 & 0.539 \\
    6  & 0.62 & 0.68 & +0.06 & 0.420 & 0.460 \\
    4  & 0.60 & 0.64 & +0.04 & 0.289 & 0.308 \\
    2  & 0.42 & 0.60 & +0.18 & 0.176 & 0.251 \\
    1  & 0.12 & 0.22 & +0.10 & 0.002 & 0.004 \\
    \bottomrule
  \end{tabular}
\end{table}

Both systems degrade as leads are removed, but the CKP-assisted condition achieves a higher
GQ at every operating point. The advantage is modest under complete information (+0.03 at
12 leads) and becomes largest in the 2-lead configuration (+0.18). The control condition falls
below the 0.60 safety threshold at 2 leads ($\mathrm{RFB} = 4$), whereas the CKP-assisted model
maintains acceptable grounding down to 2 leads ($\mathrm{RFB} = 2$), delaying the failure
boundary by two full lead levels.

\subsection{Aggregate Benchmark Measures}

Table~\ref{tab:aggregate-results} summarizes the aggregate benchmark results.

\begin{table}[H]
  \centering
  \caption{Aggregate reasoning-degradation benchmark}
  \label{tab:aggregate-results}
  \begin{tabular}{@{}lrrr@{}}
    \toprule
    \textbf{Metric} & \textbf{ECG-only} & \textbf{ECG + CKP} &
    \textbf{Improvement} \\
    \midrule
    AURDC & 0.70 & 0.81 & +0.11 absolute \\
    Relative AURDC change & --- & --- & +15.7\% \\
    RFB ($\tau=0.60$) & 4 leads & 2 leads & 2 fewer leads required \\
    GQ at 2 leads & 0.42 & 0.60 & +0.18 \\
    GQ at 1 lead & 0.12 & 0.22 & +0.10 \\
    \bottomrule
  \end{tabular}
\end{table}

The CKP condition achieves an AURDC of 0.81, exceeding the predefined success criterion of 0.80
and improving over the control AURDC of 0.70 (+15.7\% relative gain).

\section{Comparison with Existing Work}

Established ECG deep-learning benchmarks (Strodthoff et al. \cite{strodthoff2021}, Ribeiro et
al. \cite{ribeiro2020}, K-MERL \cite{kmerl2025}) evaluate models primarily on classification
accuracy under fixed lead sets. The proposed framework advances beyond prior work in three key
respects:
\begin{itemize}
  \item \textbf{Evaluation Scope:} Conventional metrics reward label accuracy regardless of
  whether cited evidence was observable. The proposed framework adds claim-level verification
  (\textsc{Supported}, \textsc{Contradicted}, \textsc{Unavailable}), detecting when correct
  labels rest on hallucinated evidence from missing leads.
  \item \textbf{Treatment of Sensor Scarcity:} Avoids generative waveform inpainting or
  synthetic reconstruction, evaluating authentic physiological degradation.
  \item \textbf{Reduced-Lead Resilience:} Demonstrates that ontological constraints preserve
  reasoning consistency under 2-lead telemetry ($GQ = 0.60$), while True Grounding ($TGR = 0.251$)
  appropriately indicates that telemetry cannot replace complete 12-lead diagnostic coverage.
\end{itemize}

\section{Research-Based Discussion and Clinical Implications}

\subsection{The Illusion of Robustness}
A central finding is that label accuracy and reasoning reliability diverge sharply as sensors are
removed. Foundation models frequently output correct disease labels by exploiting dataset
shortcuts and citing evidence in disconnected leads. The steep decline in EO and TGR exposes
this hidden failure mode, confirming that classification accuracy without physical observability
creates an illusion of clinical robustness.

\subsection{Clinical Knowledge as a Reasoning Anchor}
The SNOMED CT prior acts as an inductive constraint in the shared transformer, restricting
hypothesis attention to clinically valid finding-disease relationships. The largest gain occurs
at 2 leads (+0.18). Nevertheless, the prior does not create new physical evidence: both paradigms
collapse at 1 lead ($TGR < 0.005$), confirming that structured knowledge cannot replace missing
sensor measurements.

\subsection{Translational Significance and Limitations}
The 2-lead results confirm that ontology-guided AI can safely assist rhythm and conduction triage
in emergency telemetry. However, single-lead wearables cannot responsibly diagnose complex
regional infarctions, requiring models to abstain or recommend 12-lead acquisition. Study
limitations include reliance on one held-out partition (PTB-XL), dependency on automated
observation extractors, and the need for prospective multi-center clinical validation before
bedside deployment.

\section{Chapter Summary}
Progressive lead loss causes severe degradation in evidence observability and claim grounding.
The CKP-assisted model improves AURDC from 0.70 to 0.81 and extends the Reasoning Failure
Boundary from 4 leads down to 2 leads. These results confirm that ontological clinical priors
substantially improve reasoning resilience, while claim-level verification ensures AI safety
limits remain transparent.

\begin{thebibliography}{99}
\bibitem{ptbxl2020}
P. Wagner, N. Strodthoff, R.-D. Bousseljot, D. Kreiseler, F. I. Lunze,
W. Samek, and T. Schaeffter, "PTB-XL, a large publicly available
electrocardiography dataset," \textit{Scientific Data}, vol.~7, art.~154, 2020.
DOI: \url{https://doi.org/10.1038/s41597-020-0495-6}.

\bibitem{strodthoff2021}
N. Strodthoff, P. Wagner, T. Schaeffter, and W. Samek,
"Deep learning for ECG analysis: Benchmarks and insights from PTB-XL,"
\textit{IEEE Journal of Biomedical and Health Informatics}, vol.~25, no.~5,
pp.~1519--1528, 2021.
DOI: \url{https://doi.org/10.1109/JBHI.2020.3022989}.

\bibitem{ribeiro2020}
A. H. Ribeiro, M. H. Ribeiro, G. M. M. Paix\~ao, D. M. Oliveira,
P. R. Gomes, J. A. Canazart, M. P. S. Ferreira, C. R. Andersson,
P. W. Macfarlane, W. Meira Jr., T. B. Sch\"on, and A. L. P. Ribeiro,
"Automatic diagnosis of the 12-lead ECG using a deep neural network,"
\textit{Nature Communications}, vol.~11, art.~1760, 2020.
DOI: \url{https://doi.org/10.1038/s41467-020-15432-4}.

\bibitem{kmerl2025}
C. Liu, C. Ouyang, Z. Wan, H. Wang, W. Bai, and R. Arcucci,
"Knowledge-enhanced multimodal ECG representation learning with arbitrary-lead
inputs," in \textit{Findings of the Association for Computational Linguistics:
EMNLP 2025}, pp.~7298--7316, 2025.

\bibitem{samplepaper} A. Author and B. Author, "Title of the research paper", \textit{Journal or Conference Name}, vol.~1, no.~1, pp.~1--10, 2025.
DOI: \url{https://doi.org/10.3844/jcssp.2021.610.623}.
\bibitem{r1} Wu, F., Souza, A., Zhang, T., Fifty, C., Yu, T., and Weinberger, K, "Simplifying graph convolutional networks", \textit{In the International Conference on Machine Learning"}, pages 6861–6871. Available from \url{https://proceedings.mlr.press/v97/wu19e.html}
\bibitem{r2} Zhang, Z. and Fink, O, "Algorithm-informed graph neural networks for leakage detection and localization in water distribution networks", \textit{Reliability Engineering \& System Safety}, page 111494. DOI: \url{https://doi.org/10.1016/j.ress.2025.111494}
\bibitem{r3}, Zhao, F., Liu, J., Zhang, S., Tao, L., Wang, J., Zhang, Z., and Li, N, "A novel method for detecting the leakage duration of water pipe network based on the anchor box mechanism", \textit{In 2025 IEEE 2nd International Conference on Energy and Electrical Engineering (EEE)}, pages 1–5. IEEE, DOI: \url{https://doi.org/10.1109/eee64897.2025.11162744}
\bibitem{samplebook} C. Author, \textit{Title of the Book}, 2nd ed. City, Country: Publisher, 2024.
\bibitem{samplewebsite} Organization Name, ``Title of the webpage,'' 2026. [Online]. Available: \url{https://example.com}. [Accessed: 7-Aug-2026].
\end{thebibliography}

\end{document}
